\documentclass[11pt]{article}
\usepackage{mathunal-base}
\mathunalfuente{serif}

\renewcommand{\examtitulo}{Segundo Examen Parcial de Cálculo en Varias Variables}
\renewcommand{\examfecha}{2 de Noviembre de 2013}

\begin{document}

\examheader

\vspace{4pt}
\noindent
\begin{minipage}[t]{0.62\linewidth}
Nombre \dotfill

\vspace{4pt}
Profesor \dotfill
\end{minipage}%
\hfill
\begin{minipage}[t]{0.32\linewidth}
Doc. Identidad \dotfill

\vspace{4pt}
Grupo \dotfill
\end{minipage}

\vspace{6pt}
\noindent\textbf{Instrucciones.} El examen tiene una duración de 1 hora y 40 minutos. \textbf{No se admiten preguntas durante el examen.}

\noindent No está permitido el uso de calculadora, ni teléfonos móviles, ni cualquier otro tipo de aparato electrónico. Todos los puntos de procedimiento deben estar debidamente justificados.

\vspace{10pt}
\begin{enumerate}
\item{} (20\%) \textbf{Selección múltiple.}
\begin{enumerate}
\item{} (10\%) Considere los siguientes dos enunciados:

\vspace{4pt}
\noindent \textbf{(I)} Sea $f:\mathbb{R}^2\to\mathbb{R}$ una función de clase $C^1$. Si $(2,1)$ es un punto crítico de $f$ y $f_{xx}(2,1)f_{yy}(2,1)<[f_{xy}(2,1)]^2$, entonces $f$ tiene un punto silla en $(2,1)$.

\noindent \textbf{(II)} La integral $\displaystyle\int_0^{2\pi}\int_0^2\int_r^2 1\,dz\,dr\,d\theta$ representa el volumen encerrado por el cono $z=\sqrt{x^2+y^2}$ y el plano $z=2$.

\vspace{4pt}
\noindent Podemos afirmar que:

\vspace{4pt}
\noindent $(a)$ Los enunciados \textbf{(I)} y \textbf{(II)} son falsos.

\noindent $(b)$ Los enunciados \textbf{(I)} y \textbf{(II)} son verdaderos.

\noindent $(c)$ El enunciado \textbf{(I)} es falso y el enunciado \textbf{(II)} es verdadero.

\noindent $(d)$ El enunciado \textbf{(I)} es verdadero y el enunciado \textbf{(II)} es falso.

\item{} (10\%) El valor de la integral $\displaystyle\iiint_B e^{x+y+z}\,dV$ sobre el cubo $B=[0,1]\times[0,1]\times[0,1]$ es

\vspace{4pt}
\noindent $(a)\ (e^2-1)^3$ \hspace{1cm} $(b)\ e^3$ \hspace{1cm} $(c)\ (e+1)^3$ \hspace{1cm} $(d)\ (e-1)^3$.
\end{enumerate}

\item{} (35\%)
\begin{enumerate}
\item{} (18\%) Use multiplicadores de Lagrange para determinar cuál es el área máxima que puede tener un rectángulo si la longitud de su diagonal es 2.

\vspace{5cm}

\item{} (17\%) Evalúe la integral $\displaystyle\int_0^1\int_0^{\cos^{-1}y} e^{\operatorname{sen}x}\,dx\,dy$. Dibuje y describa la región de integración.

\vspace{5cm}
\end{enumerate}

\item{} (25\%) Sea $E$ el sólido que se encuentra por dentro de la superficie $x^2+y^2+z^2=2z$ y debajo de la superficie $z=\sqrt{3(x^2+y^2)}$. Si en cada punto $(x,y,z)\in E$ la densidad es igual a la distancia del punto al origen, halle la masa de $E$. Dibuje y describa la región de integración.

\vspace{6cm}

\item{} (20\%) Use el teorema de cambio de variables para evaluar la integral $\displaystyle\iint_D \operatorname{sen}(9x^2+4y^2)\,dA$, donde $D$ es la región del primer cuadrante acotada por la elipse $9x^2+4y^2=1$. Dibuje y describa las regiones de integración.

\vspace{6cm}
\end{enumerate}

\examfooter
\end{document}
